\chapter*{Glossary of terms, abbreviations, and acronyms}
\markboth{Glossary of terms, abbreviations, and acronyms}{Glossary of terms, abbreviations, and acronyms}

\begin{spacing}{1.5}
\begin{longtable}[l]{p{0.28\textwidth} p{0.68\textwidth}}
    \textbf{2p2h} & Two-particle-two-hole interaction resulting in two nucleons in the final state \\
        \textbf{$A$} & Mass number of a given nucleus referring to the number of nucleons (protons + neutrons) \\
        \textbf{$A$-dependence} & Nuclear dependence \\
        \textbf{ADC} & \textbf{A}nalog-to-\textbf{D}igital \textbf{C}onverter of the signal \\
        \textbf{AGKY} & \textbf{A}ndreopoulous-\textbf{G}allagher-\textbf{K}ehayias-\textbf{Y}ang model describing the hadronisation process in GENIE \\
        \textbf{Baffle} & Graphite-aluminum tube designed to partially degrade mis-steered beams to protect other components of the beamline \\
        \textbf{Calorimetric spline} & Piecewise calorimetric correction applied to reconstructed recoil energy to account for deposition in passive material, limited containment, and invisible energy deposits in the detector from particles such as neutrons    \\
        \textbf{CC} & \textbf{C}harged-\textbf{C}urrent interaction wherein (anti)neutrino exchanges a $W^{\pm}$ boson, converting into its coupled charged lepton \\
        \textbf{CP} & \textbf{C}harge-\textbf{P}arity symmetry  \\
        \textbf{DAQ} & \textbf{D}ata \textbf{A}c\textbf{Q}uisition, the system that receives raw data from the detector, digitises them, and stores them on disk \\
        \textbf{`Daisy petal'} & One of the 12 transverse bins in which the hexagonal cross-section of the scintillator tracker is divided for the flux and differential cross-section calculation \\
        \textbf{`Daisy reweight'} & Linear combination of the scintillator distributions extracted in 12 `daisy petals' to match the flux of the scintillator tracker to that of a given nuclear target in cross-section ratios  \\
        \textbf{DCNN} & \textbf{D}eep \textbf{C}onvolutional \textbf{N}eural \textbf{N}etwork used in this work to reconstruct the interaction vertex  \\
        \textbf{DIS} & \textbf{D}eep \textbf{I}nelastic \textbf{S}cattering interaction channel in which a high-energetic (anti)neutrino interacts with the constituents of nucleons, producing a hadronic shower  \\
        \textbf{Downstream} & Relative position with respect to the (anti)neutrino beam referring to a position \textit{in} the direction of the beam   \\
        \textbf{DUNE} & \textbf{D}eep \textbf{U}nderground \textbf{N}eutrino \textbf{E}xperiment, predominantly argon-based future neutrino oscillation experiment  \\
        \textbf{ECAL} & \textbf{E}lectromagnetic \textbf{CAL}orimeter designed to contain electromagnetic showers  \\
        \textbf{EMC effect} & \textbf{E}uropean \textbf{M}uon \textbf{C}ollaboration effect; one of the observed nuclear modifications to structure functions describing depletion in the bound nucleon cross-section ratio in the $0.3 \lesssim x \lesssim 0.7$   \\ 
        \textbf{FEB} & \textbf{F}ront-\textbf{E}nd \textbf{B}oard; electronics board attached on top of a PMT that digitises the output signal from the PMT \\
        \textbf{Fermi motion} & Isotropic initial state movement of nucleons inside the nuclear environment \\
        \textbf{FHC} & \textbf{F}orward \textbf{H}orn \textbf{C}urrent configuration of the magnetic horns in the beamline, resulting in neutrino-dominated beam  \\
        \textbf{FSI} & \textbf{F}inal \textbf{S}tate \textbf{I}nteractions that particles produced by (anti)neutrino interactions with nucleons inside the nucleus can undergo prior to exiting the nucleus \\
        \textbf{G4NuMI} & \textbf{G}EANT\textbf{4} package for simulation of the \textbf{NuMI} beamline  \\
        \textbf{Gate} & MINERvA's data-taking period of 16~$\mu$s   \\
        \textbf{GENIE} & \textbf{G}enerates \textbf{E}vents for \textbf{N}eutrino \textbf{I}nteraction \textbf{E}xperiments; (anti)neutrino-nucleus interaction simulation, also known as event generator, used to model (anti)neutrino interactions and make predictions  \\
        \textbf{GENIEXSecExtract} & MINERvA's cross-section validation tool, which is external to the analysis framework \\
        \textbf{GEANT4} & \textbf{GE}ometry \textbf{AN}d \textbf{T}racking software used for detector geometry and particle propagation through the detector modelling \\
        \textbf{HCAL} & \textbf{H}adronic \textbf{CAL}orimeter designed to contain hadronic showers   \\
        \textbf{Horn} & Parabolic focusing and charge-sign selecting magnet used in the (anti)neutrino beamline setup \\
        \textbf{ID} & \textbf{I}nner \textbf{D}etector  \\
        \textbf{Inclusive} & Considering all final states, i.e.~all (anti)neutrino interaction modes \\
        \textbf{INTRANUKE} & Intranuclear cascade simulation used in GENIE to model FSI  \\
        \textbf{KNO} & \textbf{K}oba-\textbf{N}ielsen-\textbf{O}lesen phenomenological model used in GENIE's AGKY hadronisation model to describe the hadronisation process for events with lower invariant mass \\
        \textbf{LE} & \textbf{L}ow \textbf{E}nergy regime of the NuMI beam and corresponding MINERvA data-taking period with (anti)neutrino flux with(anti)neutrino energy peak of $\sim$3~GeV  \\
        \textbf{LFG} & \textbf{L}ocal \textbf{F}ermi \textbf{G}as model for the initial state of the nucleons inside the nucleus taking into account the local density at the interaction vertex \\
        \textbf{MAT} & \textbf{M}INERvA \textbf{A}nalysis \textbf{T}oolkit used in MINERvA analyses mainly handling the propagation of systematic uncertainties \\
        \textbf{ME} & \textbf{M}edium \textbf{E}nergy regime of the NuMI beam and corresponding MINERvA data-taking period with (anti)neutrino flux with(anti)neutrino energy peak of $\sim$6~GeV\\
        \textbf{MEC} & \textbf{M}eson \textbf{E}xchange \textbf{C}urrent correlating nucleons inside the nuclear environment via an exchange of virtual mesons, often pions\\
        \textbf{MC} & \textbf{M}onte \textbf{C}arlo simulation refers to the (anti)neutrino interaction simulation  \\
        \textbf{Migration} & Process of smearing of the reconstructed kinematics of an event from its true kinematics due to finite detector resolution and imperfect reconstruction described by migration matrix  \\
        \textbf{MINERvA} & \textbf{M}ain \textbf{I}njector \textbf{N}eutrino \textbf{E}xpe\textbf{R}iment to study \textbf{$\nu$}--\textbf{A} interactions, a dedicated cross-section experiment  \\
        \textbf{MINOS} & \textbf{M}ain \textbf{I}njector \textbf{N}eutrino \textbf{O}scillation \textbf{S}earch, a neutrino oscillation experiment; in the context of this work often refers to its near detector only, which was positioned downstream of the MINERvA detector  \\
        \textbf{ML} & \textbf{M}achine \textbf{L}earning, refers to the technique used to reconstruct the interaction vertex \\
        \textbf{Module} & MINERvA's detector tracking unit which refers to two planes of scintillator strips  \\
        \textbf{MIP} & \textbf{M}inimum \textbf{I}onising \textbf{P}article such as (anti)muon   \\
        \textbf{NC} & \textbf{N}eutral-\textbf{C}urrent interaction wherein (anti)neutrino exchanges a $Z$ boson with no charge transfer and leaves the interaction without changing to its paired charged lepton \\
        \textbf{NEUT} & Another (anti)neutrino event generator predominantly used in Japanese experiments, such as the neutrino oscillation experiment T2K  \\
        \textbf{Nuclear target} & Passive material interlaid between the active scintillator planes in MINERvA consisting of combinations of carbon, iron, and lead in the context of this work (there is also a water and helium target)  \\
        \textbf{NuMI} & \textbf{N}eutrinos ({$\bm{\nu}$}) at the \textbf{Main} \textbf{I}njection; (anti)neutrino beam production setup at Fermilab   \\
        \textbf{OD} & \textbf{O}uter \textbf{D}etector   \\
        \textbf{PDFs} & \textbf{P}arton \textbf{D}istribution \textbf{F}unctions describing the distribution of partons within the nucleus  \\
        \textbf{PMNS matrix} & \textbf{P}ontecorvo-\textbf{M}aki-\textbf{N}akagawa-\textbf{S}akata mixing matrix relating the neutrino mass eigenstates and the neutrino flavour eigenstates, describing neutrino oscillation  \\
        \textbf{PMT} & \textbf{P}hoto\textbf{M}ultiplier \textbf{T}ube with 64 channels converting the collected light into charge \\
        \textbf{Plane} & MINERvA's detector tracking unit made of scintillator strips, with either $X$, $U$, or $V$ orientation of the scintillator strips (parallel to the $y$-axis, 60$^{\circ}$ clockwise and anticlockwise, respectively)   \\
        \textbf{Plastic} & Refers to the hydrocarbon scintillator used to make the active tracker region or the scintillator planes in between the nuclear targets  \\
        \textbf{Playlist} & Lists of MINERvA runs and subruns that correspond to a specific detector configuration during a data-taking period \\
        \textbf{POT} & \textbf{P}rotons \textbf{O}n \textbf{T}arget; a metric used to characterise the integrated luminosity of the (anti)neutrino beam \\
        \textbf{PPFX} & \textbf{P}ackage to \textbf{P}redict the \textbf{F}lu\textbf{X} used to constrain flux predictions by external data and evaluate systematic uncertainties  \\
        \textbf{PYTHIA/JETSET} & Standard tool for generation of high-energy collisions, spanning models for few-body hard processes to complex multihadronic final states used in GENIE's AGKY hadronisation model to describe the hadronisation process for events with higher invariant mass \\
        \textbf{Recoil} & Refers to either the total outgoing hadronic system in an interaction, or more specifically to its energy \\
        \textbf{RES} & \textbf{RES}onant pion production interaction channel in which the (anti)neutrino is energetic enough to excite a nucleon into a baryon resonance which then decays, most often to pions \\
        \textbf{RFG} & \textbf{R}elativistic \textbf{F}ermi \textbf{G}as model for initial state of the nucleons inside the nucleus, treating nucleons as a collection of non-interacting fermions in a constant nuclear potential \\
        \textbf{RHC} & \textbf{R}everse \textbf{H}orn \textbf{C}urrent configuration of the magnetic horns in the beamline, resulting in antineutrino-dominated beam   \\
        \textbf{RPA} & \textbf{R}andom \textbf{P}hase \textbf{A}pproximation accounting for long-range correlations in (anti)neutrino-nucleus scattering, arising from the polarisation of the nuclear medium by the electroweak propagator \\
        \textbf{SF} & \textbf{S}pectral \textbf{F}unction model for the initial state of the nucleons inside the nucleus, taking into account two- and three-body nucleon-nucleon interactions, thereby modifying the available nuclear quantum states \\
        \textbf{Shadowing} & Suppression of the cross-section of bound nucleons due to destructive interference of hadronic fluctuations of the intermediate boson that multiple-scatter off surface nucleons at $x \lesssim 0.1$ compared to a free nucleon  \\
        \textbf{SIS} & \textbf{S}hallow \textbf{I}nelastic \textbf{S}cattering referring to the transition region between RES and DIS interaction channels with non-resonant pion production and multi-pion final states, defined by \mbox{$(m_N + m_{\pi}) < W < 2$ GeV/$c^2$} and $Q^2$ around 1 (GeV/$c$)$^2$ \\
        \textbf{SM} & \textbf{S}tandard \textbf{M}odel of particle physics, categorising all elementary particles and describing three out of the four known fundamental forces (excluding gravity)  \\
        \textbf{Smearing} & Describes the process that results in a difference between the actual kinematics of the physics that occurred, and the kinematics of the physics measured (reconstructed) in the detector due to finite detector resolution and imperfect reconstruction    \\
        \textbf{SRC} & \textbf{S}hort-\textbf{R}ange \textbf{C}orrelations between nucleons arising from brief overlaps of their wavefunctions  \\
        \textbf{Strips} & Triangular scintillator bars which make up one scintillator plane, the basic MINERvA tracking unit  \\
        \textbf{SuSA} & Super-scaling model based on phenomenon shown by electron-nucleus scattering data, in the context of this wok refers to an alternative 2p2h model \\
        \textbf{TDC} & \textbf{T}ime-to-\textbf{D}igital \textbf{C}onverter of the signal timing  \\
        \textbf{Time slice} & Time window with a variable length within the gate in which reconstruction steps are performed  \\
        \textbf{true/truth} & Direct output of the simulated prediction of the (anti)neutrino event generator (no reconstruction) \\
        \textbf{Unfolding} & Process of reverting the smearing of reconstructed distribution and estimating the true distribution from the measured smeared spectrum in order to report detector-independent cross-section  \\
        \textbf{Upstream} &  Relative position with respect to the (anti)neutrino beam referring to a position \textit{against} the direction of the beam \\
        \textbf{Valencia model} & Refers to a model used to describe 2p2h and RPA \\
        \textbf{Warping study} & Unfolding study using simulations that have been reweighted to depict alternative physics models, exaggerate features of the underlying distributions, or capture differences in the shapes of distributions between the data and the original underlying simulation to evaluate model dependence in unfolding \\
        \textbf{WLS} & \textbf{W}ave\textbf{L}ength-\textbf{S}hifting; usually refers to a fibre that can change the wavelength of the collected light \\
\end{longtable}
\end{spacing}
